\section{Related Work}
\label{sec:related}

% Sources: RELATED_WORK.md (verification record and allowed sentences). Depth is metadata and abstract only; claims about a work are "per its abstract" (X8: no performance comparison was run).
Table~\ref{tab:related} maps the verified set.
Rows 1--10 are the works of the project's reference list; rows 11--13 were filled from their abstracts only (verified 2026-10-15).
A cross means that the abstract does not describe the property, not that the paper lacks it.
No performance comparison with any of these works was run.

\begin{table*}[t]
\centering
% src: results/s10/tables/related_work.csv (rows copied; row 'this' wording as in the file)
\caption{Prior art, mapped. Check: described; $\sim$: partly or in a different sense; --: not described in the abstract. The structure-aware mark for this work records that message passing is implemented, not that it helped.}
\label{tab:related}
\small
\begin{tabular}{@{}rp{5.4cm}cp{3.0cm}cc@{}}
\toprule
Ref & Work & Learned trust & Structure-aware & Vehicular offloading & Explainable \\
\midrule
1 & GNN-Mamba spatiotemporal trust evaluation & \checkmark & \checkmark & -- & -- \\
2 & Multi-hop trust via GNN-aided agentic AI & \checkmark & \checkmark & -- & -- \\
3 & Federated GNN multi-agent RL for AoI in VEC & -- & \checkmark & \checkmark & -- \\
4 & Joint allocation under asymmetric information & $\sim$ (economic) & -- & \checkmark & -- \\
5 & DRL for delay-optimized VFC offloading & -- & -- & \checkmark & -- \\
6 & Explainability-as-a-Service for edge AI & -- & -- & -- & \checkmark \\
7 & XRL for RAN slicing (6G ORAN) & -- & -- & -- & \checkmark \\
8 & Self-explaining RL for mobile resource allocation & -- & -- & -- & \checkmark \\
9 & XAI for RL-based networking & -- & -- & -- & \checkmark \\
10 & Explainable fraud detection (GNNExplainer + Shapley) & -- & \checkmark & -- & \checkmark \\
11 & SGA-SAC: Bayesian Stackelberg + graph-attention DRL, explainable VEC offloading (Cheng et al., 2026) & -- & \checkmark (graph attention) & \checkmark & \checkmark (Grad-CAM) \\
12 & Trust-based malicious RSU detection in edge-enabled VANETs (Siddiqua and Jahan) & -- & $\sim$ (RSU-RSU interaction) & -- & -- \\
13 & G-TRAC: trust-aware routing for distributed GenAI inference at the edge (Nguyen and Elmroth, 2026) & -- & $\sim$ (path search over peers) & -- & -- \\
\midrule
this & This work & \checkmark & \checkmark (built; no measured gain over node-local) & \checkmark & \checkmark \\
\bottomrule
\end{tabular}
\end{table*}

\paragraph*{Learned, structure-aware trust}
% src: RELATED_WORK.md, positioning sentence 1
References~\cite{r1,r2} already do learned, structure-aware trust with GNNs; they share authorship and are one continuing line of work.
We adopt the formulation and integrate it into the offloading objective, and our own result is that the graph did not help in our scenario (Sec.~\ref{sec:graph}).

\paragraph*{Offloading in vehicular edge and fog computing}
References~\cite{r3,r4,r5} address offloading or resource allocation in vehicular edge and fog computing: a federated GNN multi-agent reinforcement-learning method for age of information~\cite{r3}, joint allocation under asymmetric information~\cite{r4}, and deep reinforcement learning for delay-optimised offloading~\cite{r5}.
By the table, only~\cite{r3} is structure-aware, and none of the three is marked explainable.

\paragraph*{Explainability}
References~\cite{r6,r7,r8,r9} concern explainability for edge AI, RAN slicing, mobile-network resource allocation and RL-based networking; reference~\cite{r10} combines GNNExplainer with Shapley values for fraud detection.
% src: RELATED_WORK.md table (titles of rows 6-10); no statement beyond the table
Reference~\cite{r11}, SGA-SAC, is the closest explainable vehicular-offloading work we found.
Per its abstract, it models vehicles and RSUs as a hierarchical Bayesian Stackelberg game, uses soft actor-critic with a graph attention mechanism in the policy network, and uses Grad-CAM to visualise the multi-agent decision, evaluated on latency and energy consumption; the abstract does not describe a node-trust model.
% src: RELATED_WORK.md item 11 (allowed differentiation)
Its explanation is a gradient visualisation of a learned policy, whereas ours is an exact decomposition of a fixed linear score plus a separate attribution on the trust head, and its objective is latency and energy whereas ours is deadline success under injected node degradation.
The contribution difference is the exactly decomposable score and a trust head under injected degradation, not a performance claim.

\paragraph*{Behavioural trust of RSUs, and trust floors}
Reference~\cite{r12} is prior art for behavioural trust of RSUs: per its abstract, a trust-based scheme that uses stable, frequent RSU-to-RSU interaction and a check for altered sensor-detected data in beacons to detect malicious RSUs in edge-enabled VANETs, with a self-reported detection of about 92\% of malicious RSUs; the abstract does not describe offloading, a learned model or an explanation.
% src: RELATED_WORK.md item 12
Reference~\cite{r13} formulates routing of distributed generative-AI inference over edge devices as a risk-bounded shortest-path problem solved by trust-floor pruning combined with Dijkstra's search; it is not vehicular and concerns routing of inference, not task offloading from vehicles.
Its trust floor, which removes candidates below a threshold before selection, is the nearest analogue of the gate we test as a baseline (Gate-C, Sec.~\ref{sec:method}); in our scenario the gate and the weighted rule performed equivalently (Sec.~\ref{sec:gate}).
Gate-C is our own construction on our own trust head; we did not implement~\cite{r13}.
